% year: תשע"ח
% type: מבחן
% semester: ב
% moed: ב
% number: 4א
% points: 15
% categories: continuity, func-limits
% source: exams/מבחנים/תשעח/מועד ב תשעח סמסטר ב.pdf

\begin{question}
לאילו ערכים של של הפרמטרים $a$ ו-$b$ הפונקציה
\[
f(x)=\begin{cases}\dfrac{\tan(2x)-x}{x+\sin(2x)} & -\frac\pi4<x<0\\[1ex] ax+b & 0\le x\le1\\[0.5ex] x^{1/(x-1)} & x>1\end{cases}
\]
תהיה רציפה בתחום $x>-\frac\pi4$?
\end{question}

\begin{hint}
הפונקציה רציפה בתוך כל אחד משלושת הקטעים; יש לבדוק רק את נקודות התפר $x=0$ ו-$x=1$. ב-$x=0$ חלקו מונה ומכנה ב-$x$; ב-$x=1$ זהו גבול מהצורה $1^\infty$.
\end{hint}

\begin{solution}
\textbf{רציפות בתוך הקטעים.} ב-$\left(-\frac\pi4,0\right)$: $\tan2x$ מוגדרת ורציפה (כי $2x\in\left(-\frac\pi2,0\right)$), והמכנה $x+\sin2x$ שלילי ממש (שני המחוברים שליליים), ולכן $f$ רציפה שם כמנה של רציפות. ב-$(0,1)$ $f$ פולינום. ב-$(1,\infty)$: $x^{1/(x-1)}=e^{\frac{\ln x}{x-1}}$ רציפה. נותר לבדוק את $x=0$ ו-$x=1$.

\textbf{הנקודה $x=0$.} $f(0)=b$, והגבול מימין הוא $b$. משמאל, נחלק מונה ומכנה ב-$x$:
\[
\lim_{x\to0^-}\frac{\tan2x-x}{x+\sin2x}=\lim_{x\to0^-}\frac{2\cdot\frac{\tan2x}{2x}-1}{1+2\cdot\frac{\sin2x}{2x}}=\frac{2-1}{1+2}=\frac13,
\]
לפי הגבולות היסודיים $\frac{\sin t}{t}\to1$, $\frac{\tan t}{t}\to1$. לכן צריך $b=\frac13$.

\textbf{הנקודה $x=1$.} $f(1)=a+b$, והגבול משמאל הוא $a+b$. מימין:
\[
\lim_{x\to1^+}x^{1/(x-1)}=\lim_{x\to1^+}e^{\frac{\ln x}{x-1}}=e^1=e,
\]
כי $\lim_{x\to1}\frac{\ln x}{x-1}=1$ (למשל לפי לופיטל: $\frac{1/x}{1}\to1$, או לפי $\ln(1+t)/t\to1$ עם $t=x-1$), ולפי רציפות האקספוננט. לכן צריך $a+b=e$, כלומר $a=e-\frac13$.

\textbf{תשובה:} $a=e-\frac13$, $b=\frac13$.
\end{solution}
